\documentclass[a4paper, 11pt]{ctexart}

% 页面边距与排版宏包
\usepackage{geometry}
\geometry{left=2.5cm, right=2.5cm, top=2.5cm, bottom=2.5cm}
\usepackage{enumitem}
\usepackage{xcolor}
\usepackage{ulem}

% 自定义标题格式
\title{\textbf{\Huge 深蓝国际象棋协会}\\ \vspace{0.5em} \Large “薪火传承”校园挑战赛题库}
\author{}
\date{2026年4月18日}

\begin{document}

\maketitle

\newpage
% ==========================================
% 第一部分：纯题目版本（供打印发放使用）
% ==========================================
\section*{【纯题目版】}
\textit{欢迎来到深蓝国际象棋协会“薪火传承”活动现场！快来测试一下你的“棋商”与科大社团底蕴吧。}
\vspace{1em}

\begin{enumerate}[label=\arabic*., itemsep=1.5em]

    \item \textbf{创造历史！中国第一位夺得国际象棋女子世界冠军的棋手是谁？} \\
    A. 侯逸凡 \qquad B. 谢军 \qquad C. 居文君 \qquad D. 诸宸

    \item \textbf{中国科大校学生深蓝国际象棋协会正式成立于哪一年？} \\
    A. 1997年 \qquad B. 2003年 \qquad C. 2014年 \qquad D. 2020年

    \item \textbf{深蓝协会秉承的社团理念是什么？} \\
    A. 友谊第一，比赛第二 \qquad B. 棋盘演绎风云变幻，布阵提升棋艺思维 \\
    C. 探索星辰大海 \qquad \qquad \, D. 培养特级大师

    \item \textbf{中国科大校园内的首届“深蓝杯”国际象棋比赛是在哪一天诞生的？} \\
    A. 4月18日 \qquad B. 10月4日 \qquad C. 11月11日 \qquad D. 1月1日

    \item \textbf{2023年，哪位中国棋手迎难而上，成为中国首位夺得国际象棋男子个人世界冠军的棋手，补齐了中国国象“四步走”战略的最后一块拼图？} \\
    A. 韦奕 \qquad B. 王玥 \qquad C. 丁立人 \qquad D. 余泱漪

    \item \textbf{在深蓝协会热推的“原子棋 (Atomic)”规则中，当棋子被吃掉发生爆炸时，以下哪种棋子是“防爆”的（不会因为邻近爆炸而消失）？} \\
    A. 王 (King) \qquad B. 王后 (Queen) \qquad C. 兵 (Pawn) \qquad D. 马 (Knight)

    \item \textbf{咱们深蓝协会在今年 4 月举办的主题趣味赛，是致敬了哪部著名的国际象棋影视剧？} \\
    A. 《少年班》 \qquad B. 《生活大爆炸》 \qquad C. 《后翼弃兵》 \qquad D. 《哈利波特》

    \item \textbf{在深蓝协会推广的“疯狂乱斗 (Crazyhouse)”变体中，吃掉的棋子变色后可以放回棋盘，这种操作术语叫什么？} \\
    A. 召唤 \qquad B. 空降 (Drop) \qquad C. 闪现 \qquad D. 投胎

    \item \textbf{在国际象棋最高等级分排名中，中国曾有一位男子棋手历史性地闯入了“2800分俱乐部”，他是谁？} \\
    A. 丁立人 \qquad B. 韦奕 \qquad C. 王玥 \qquad D. 卜祥志

    \item \textbf{在“将军三次即获胜 (3-Check)”变体中，如果一次移动同时达成“双反击将军”，算几次将军？} \\
    A. 1次 \qquad B. 2次 \qquad C. 3次 \qquad D. 0次

\item \textbf{中国科大校学生深蓝国际象棋协会正式成立于哪一年？} \\
    A. 1997年 \qquad B. 2003年 \qquad C. 2010年 \qquad D. 2014年

    \item \textbf{深蓝协会秉承的宗旨是：“\uline{\hspace{3em}}演绎风云变幻，\uline{\hspace{3em}}提升棋艺思维”。} \\
    A. 棋局/对垒 \qquad B. 棋盘/布阵 \qquad C. 战场/谋略 \qquad D. 智力/计算

    \item \textbf{2003年10月4日，深蓝协会在校内创办了哪项极具品牌影响力的系列赛事？} \\
    A. 薪火杯 \qquad B. 科技杯 \qquad C. 深蓝杯 \qquad D. 巅峰杯

    \item \textbf{在明日（4月19日）的跨校友谊赛中，深蓝协会将与哪所高校的棋队展开线上对决？} \\
    A. 清华大学 \qquad B. 北京大学 \qquad C. 浙江大学 \qquad D. 复旦大学

    \item \textbf{现代国际象棋的规则（特别是王后和象的强力走法）大约是在哪个时期定型的？} \\
    A. 公元6世纪的印度 \quad B. 15世纪末的欧洲 \quad C. 19世纪的俄国 \quad D. 20世纪的美国

    \item \textbf{历史上第一位官方公认的国际象棋世界冠军是？（1886年夺冠）} \\
    A. 保罗·摩菲 \qquad B. 威廉·斯坦尼茨 \qquad C. 卡帕布兰卡 \qquad D. 菲舍尔

    \item \textbf{国际象棋联合会（FIDE）于1924年成立于哪个城市？} \\
    A. 伦敦 \qquad B. 纽约 \qquad C. 巴黎 \qquad D. 莫斯科

    \item \textbf{中国第一位夺得国际象棋女子世界冠军，打破欧洲多年垄断的棋手是谁？} \\
    A. 诸宸 \qquad B. 谢军 \qquad C. 许昱华 \qquad D. 侯逸凡

    \item \textbf{1956年，中国将国际象棋列为正式体育项目，当时大力推动此决策的元帅是？} \\
    A. 朱德 \qquad B. 彭德怀 \qquad C. 陈毅 \qquad D. 贺龙

    \item \textbf{浙江哪座城市因为名手辈出，被授予中国首个“国际象棋之城”的美誉？} \\
    A. 杭州 \qquad B. 宁波 \qquad C. 温州 \qquad D. 绍兴

    \item \textbf{中国天才少女侯逸凡在多少岁时首次夺魁，成为历史上最年轻的棋后？} \\
    A. 14岁 \qquad B. 16岁 \qquad C. 18岁 \qquad D. 20岁

    \item \textbf{在原子棋（Atomic）规则中，如果你用王直接吃掉对方的棋子，会发生什么？} \\
    A. 你赢了 \qquad B. 你的王会自爆，你输了 \qquad C. 对方的王消失 \qquad D. 无事发生

    \item \textbf{在“疯狂乱斗（Crazyhouse）”中，吃掉的棋子可以空降回棋盘，但“兵”不能空降在？} \\
    A. 1行和8行 \qquad B. 2行和7行 \qquad C. 4行和5行 \qquad D. 任何行都可以

    \item \textbf{国际象棋历史上，统治世界冠军宝座时间最长的棋手是？（长达27年）} \\
    A. 拉斯克 \qquad B. 卡斯帕罗夫 \qquad C. 卡尔森 \qquad D. 菲舍尔

    \item \textbf{1997年击败世界冠军卡斯帕罗夫的超级电脑“深蓝”是由哪家公司研发的？} \\
    A. Google \qquad B. IBM \qquad C. Microsoft \qquad D. Apple

    \item \textbf{深蓝协会的另一项招牌赛事“闪电杯”超快棋赛，首届比赛是在哪一年举行的？} \\
    A. 2020年 \qquad B. 2021年 \qquad C. 2022年 \qquad D. 2023年

    \item \textbf{在第二届“闪电杯”中，哪位同学以8战全胜的统治级表现封王，创造了赛事史上首个全胜夺冠纪录？} \\
    A. 张想亮 \qquad B. 高鹤宇 \qquad C. 陈睿 \qquad D. 叶天骐

    \item \textbf{关于明日（4月19日）与浙大的友谊赛，专业组第一阶段采用的赛制模式是？} \\
    A. 瑞士轮 \qquad B. 循环赛 \qquad C. 团队竞技场 (Team Battle Arena) \qquad D. 淘汰赛

    \item \textbf{在浙大 x 中科大友谊赛的决赛阶段，每方选手的对局用时是多少？} \\
    A. 3+2 \qquad B. 5+2 \qquad C. 10+5 \qquad D. 15+10

    \item \textbf{国际象棋中，如果一方的王没有被将军，但己方所有棋子都没有合法走法，这被称为什么？} \\
    A. 将死 (Checkmate) \qquad B. 逼和 (Stalemate) \qquad C. 弃权 \qquad D. 超时

    \item \textbf{在“双人蝉联象棋 (Bughouse)”变体中，最核心的规则特点是什么？} \\
    A. 棋子会爆炸 \qquad B. 两人组队且队友吃掉的棋子可以交给你使用 \qquad C. 王可以走两步 \qquad D. 没有王后

    \item \textbf{国际棋联 (FIDE) 授予的最高永久头衔是什么？} \\
    A. 国际大师 (IM) \qquad B. 棋协大师 (CM) \qquad C. 特级大师 (GM) \qquad D. 候补大师 (NM)

    \item \textbf{当一名兵移动到对方底线时，它可以升变为除了哪种棋子以外的任何兵种？} \\
    A. 王后 \qquad B. 王 \qquad C. 车 \qquad D. 马

    \item \textbf{在“吃过路兵”规则中，这个动作必须在对方兵前进两格后的第几步内完成？} \\
    A. 立即下一步 \qquad B. 两步内 \qquad C. 任何时候都可以 \qquad D. 对局结束前

    \item \textbf{深蓝协会目前正致力于建立数字化档案，主要整理的内容包括哪些？} \\
    A. 2003年建社以来的照片 \quad B. 荣誉证书 \quad C. 历届杯赛棋谱 \quad D. 以上皆是

    \item \textbf{在国际象棋中，所谓的“兵的升变”，兵不能升变为以下哪种棋子？} \\
    A. 王后 \qquad B. 车 \qquad C. 国王 \qquad D. 骑士（马）

    \item \textbf{关于“王车易位”，以下哪种情况依然可以进行易位？} \\
    A. 国王正在被“将军” \qquad B. 易位时国王经过的格子受对方攻击 \\
    C. 王和车之间有其他棋子 \qquad D. 参与易位的“车”受到攻击

    \item \textbf{在常规的棋子分值估算中，一只“象”通常被认为价值多少个“兵”？} \\
    A. 1个 \qquad B. 3个 \qquad C. 5个 \qquad D. 9个

    \item \textbf{“吃过路兵”动作必须在对方兵双格前冲后的第几回合立即完成？} \\
    A. 立即下一回合 \qquad B. 两个回合内 \qquad C. 三个回合内 \qquad D. 任何时候均可

    \item \textbf{世界著名的开局“西西里防御”，其标志性的第一步黑方走法是？} \\
    A. 1... e5 \qquad B. 1... c5 \qquad C. 1... d5 \qquad D. 1... Nf6

    \item \textbf{在开局阶段，以下哪项原则通常不被推荐？} \\
    A. 尽快出动轻子（马、象） \qquad B. 尽早控制中心 \\
    C. 尽早进行王车易位 \qquad D. 在没有明确目标时多次移动同一个棋子

    \item \textbf{“四步将杀（学子棋）”主要是利用了黑方哪一个格子的防御薄弱？} \\
    A. f7 \qquad B. e7 \qquad C. d7 \qquad D. g7

    \item \textbf{以下哪种开局属于“封闭式开局”？} \\
    A. 意大利开局 \qquad B. 王翼弃兵 \qquad C. 后翼弃兵 \qquad D. 西班牙开局

    \item \textbf{在单王对单王的残局中，一方想要保护兵升变，需要占据所谓的“\uline{\hspace{2em}}”？} \\
    A. 侧翼 \qquad B. 对王（Opposition） \qquad C. 边线 \qquad D. 禁区

    \item \textbf{“正方形法则”是残局中用来快速计算什么的工具？} \\
    A. 王能否追上对方的兵 \qquad B. 象能否控制长斜线 \\
    C. 两个兵如何防守 \qquad D. 战术组合的深度

    \item \textbf{当一方被将军，但没有任何合法的走法来摆脱，这被称为？} \\
    A. 逼和 \qquad B. 抽吃 \qquad C. 捉双 \qquad D. 将死

    \item \textbf{一个“马”和一个“象”配合，理论上能否将死对方的单王？} \\
    A. 绝对不能 \qquad B. 只有在对方配合时能 \qquad C. 理论上可以，但需要极高技巧 \qquad D. 算和棋

    \item \textbf{国际象棋公认的最早雏形“恰图兰卡（Chaturanga）”起源于哪个古文明？} \\
    A. 古中国 \qquad B. 古埃及 \qquad C. 古印度 \qquad D. 古波斯

    \item \textbf{被称为“国际象棋史上的莫扎特”，并在1972年赢得“世纪之战”的美国棋手是？} \\
    A. 鲍比·菲舍尔 \qquad B. 保罗·摩菲 \qquad C. 卡斯帕罗夫 \qquad D. 卡尔森

    \item \textbf{1997年，IBM公司的哪一台超级电脑第二次挑战并击败了当时的棋王卡斯帕罗夫？} \\
    A. 阿法狗 (AlphaGo) \qquad B. 深蓝 (Deep Blue) \qquad C. 沃森 (Watson) \qquad D. 曼哈顿

\item \textbf{在国际象棋对局中，如果出现“三次重复局面”，裁判通常会如何判决？} \\
    A. 白方胜 \qquad B. 黑方胜 \qquad C. 和棋 \qquad D. 双方各扣一分

    \item \textbf{关于“长距离王车易位（大易位）”，国王最终会落在哪个格子上？} \\
    A. c1或c8 \qquad B. b1或b8 \qquad C. d1或d8 \qquad D. g1或g8

    \item \textbf{下列哪种情况下，一方可以合法地宣告对局为“50步和棋”？} \\
    A. 连续50步内没有棋子被吃 \quad B. 连续50步内没有动过兵 \\
    C. 连续50步内既没有棋子被吃，也没有动过兵 \quad D. 总回合数达到50回合

    \item \textbf{如果一方的国王被将军，但唯一的摆脱方式是“吃过路兵”，这是否合法？} \\
    A. 合法 \qquad B. 不合法 \qquad C. 视具体比赛规则而定 \qquad D. 只能在快棋中执行

    \item \textbf{“意大利开局（Italian Game）”的标志性走法是白方将象移动到哪个位置？} \\
    A. b5 \qquad B. c4 \qquad C. d3 \qquad D. e2

    \item \textbf{在“法兰西防御（French Defence）”中，黑方通常在第一步如何回应白方的 1.e4？} \\
    A. 1... e6 \qquad B. 1... c6 \qquad C. 1... Nf6 \qquad D. 1... d5

    \item \textbf{被称为“后翼弃兵”的开局，白方牺牲 c4 位置的兵，核心目的是为了争夺什么的控制权？} \\
    A. 王翼 \qquad B. 边线 \qquad C. 中心 \qquad D. 底线

    \item \textbf{在开局阶段，如果一方过早地把“后”派到战场中心，最容易遭受的风险是？} \\
    A. 被对方马、象等轻子通过攻击“后”来赚取步数（时机） \qquad B. 导致无法王车易位 \\
    C. 自己的兵阵会变弱 \qquad D. 直接被对方将死

    \item \textbf{在单兵残局中，如果对方国王成功挡在了你的兵前行路径的正前方，这种防守策略常被称为？} \\
    A. 封锁 \qquad B. 对王 \qquad C. 牵制 \qquad D. 抽吃

    \item \textbf{在残局“卢塞纳位置（Lucena Position）”中，优势方通过什么技巧来保护兵升变？} \\
    A. 搭桥 \qquad B. 闷杀 \qquad C. 闪击 \qquad D. 弃子

    \item \textbf{“菲利多尔位置（Philidor Position）”是研究哪类残局时最著名的防守范式？} \\
    A. 后对兵 \qquad B. 单象对单马 \qquad C. 车兵残局 \qquad D. 双马对单王

    \item \textbf{战术“捉双（Fork）”最容易由哪种具有“跳跃性”走法的棋子完成？} \\
    A. 象 \qquad B. 车 \qquad C. 马 \qquad D. 后

    \item \textbf{世界国际象棋联合会（FIDE）的官方总部目前设在哪个国家？} \\
    A. 法国 \qquad B. 俄罗斯 \qquad C. 瑞士 \qquad D. 阿联酋

    \item \textbf{20世纪初，哪位古巴棋手因其极其精准的残局技术被誉为“国际象棋机器”？} \\
    A. 拉斯克 \qquad B. 卡帕布兰卡 \qquad C. 阿廖欣 \qquad D. 塔尔

    \item \textbf{在1996年和1997年的人机大战中，超级电脑“深蓝”挑战的世界冠军是哪位著名的“棋坛常青树”？} \\
    A. 卡斯帕罗夫 \qquad B. 卡尔波夫 \qquad C. 阿南德 \qquad D. 克拉姆尼克

    \item \textbf{目前（2026年）国际象棋等级分的世界纪录保持者（最高曾达2882分）是谁？} \\
    A. 鲍比·菲舍尔 \qquad B. 马格努斯·卡尔森 \qquad C. 丁立人 \qquad D. 涅波姆尼亚奇

\end{enumerate}

\newpage

% ==========================================
% 第二部分：带解析版本（供摊位工作人员核对使用）
% ==========================================
\section*{【带解析版：工作人员内部使用】}
\textit{注：红色为正确答案，蓝色字体为现场互动科普解说词。}
\vspace{1em}

\begin{enumerate}[label=\arabic*., itemsep=1.5em]

    \item \textbf{创造历史！中国第一位夺得国际象棋女子世界冠军的棋手是谁？} \\
    A. 侯逸凡 \qquad \textcolor{red}{\textbf{B. 谢军}} \qquad C. 居文君 \qquad D. 诸宸 \\
    \textcolor{blue}{\textbf{【现场解说】}：1991年，谢军打破了欧洲对女子世界冠军长达64年的垄断，这可以说是中国国象“薪火”点燃的标志性事件！}

    \item \textbf{中国科大校学生深蓝国际象棋协会正式成立于哪一年？} \\
    A. 1997年 \qquad \textcolor{red}{\textbf{B. 2003年}} \qquad C. 2014年 \qquad D. 2020年 \\
    \textcolor{blue}{\textbf{【现场解说】}：咱们深蓝协会建社于2003年！从创立之初，我们就致力于在校园内普及这项富有策略与美感的智力运动。}

    \item \textbf{深蓝协会秉承的社团理念是什么？} \\
    A. 友谊第一，比赛第二 \qquad \textcolor{red}{\textbf{B. 棋盘演绎风云变幻，布阵提升棋艺思维}} \\
    C. 探索星辰大海 \qquad \qquad \, D. 培养特级大师 \\
    \textcolor{blue}{\textbf{【现场解说】}：这是写在咱们社团历史里的原话！在这里，大家不仅切磋棋艺，更是锻炼沉着心态和严密的逻辑计算能力。}

    \item \textbf{中国科大校园内的首届“深蓝杯”国际象棋比赛是在哪一天诞生的？} \\
    A. 4月18日 \qquad \textcolor{red}{\textbf{B. 10月4日}} \qquad C. 11月11日 \qquad D. 1月1日 \\
    \textcolor{blue}{\textbf{【现场解说】}：2003年10月4日，首届“深蓝杯”诞生。这杯赛承载着无数科大国象爱好者的青春回忆，也是咱们社团的传统招牌！}

    \item \textbf{2023年，哪位中国棋手迎难而上，成为中国首位夺得国际象棋男子个人世界冠军的棋手，补齐了中国国象“四步走”战略的最后一块拼图？} \\
    A. 韦奕 \qquad B. 王玥 \qquad \textcolor{red}{\textbf{C. 丁立人}} \qquad D. 余泱漪 \\
    \textcolor{blue}{\textbf{【现场解说】}：丁立人在经历了惊心动魄的快棋加赛后夺冠，这让中国同时包揽了男女个人世界冠军，补齐了最后一块拼图！}

    \item \textbf{在深蓝协会热推的“原子棋 (Atomic)”规则中，当棋子被吃掉发生爆炸时，以下哪种棋子是“防爆”的？} \\
    A. 王 (King) \qquad B. 王后 (Queen) \qquad \textcolor{red}{\textbf{C. 兵 (Pawn)}} \qquad D. 马 (Knight) \\
    \textcolor{blue}{\textbf{【现场解说】}：在原子棋的爆炸波及范围内，兵是唯一能存活下来的棋子（除非直接被吃掉）。没玩过？赶紧现场来摊位对局体验一下爆炸的艺术！}

    \item \textbf{咱们深蓝协会在今年 4 月举办的主题趣味赛，是致敬了哪部著名的国际象棋影视剧？} \\
    A. 《少年班》 \qquad B. 《生活大爆炸》 \qquad \textcolor{red}{\textbf{C. 《后翼弃兵》}} \qquad D. 《哈利波特》 \\
    \textcolor{blue}{\textbf{【现场解说】}：不仅有烧脑的变体，咱们本月正在举办致敬《后翼弃兵》的趣味赛，让你体验一回主角般的下棋爽感！}

    \item \textbf{在深蓝协会推广的“疯狂乱斗 (Crazyhouse)”变体中，吃掉的棋子变色后可以放回棋盘，这种操作术语叫什么？} \\
    A. 召唤 \qquad \textcolor{red}{\textbf{B. 空降 (Drop)}} \qquad C. 闪现 \qquad D. 投胎 \\
    \textcolor{blue}{\textbf{【现场解说】}：Crazyhouse 的核心魅力就是“空降”！把你吃掉的敌方单位变成自己的援军，局势随时可能两级反转。}

    \item \textbf{在国际象棋最高等级分排名中，中国曾有一位男子棋手历史性地闯入了“2800分俱乐部”，他是谁？} \\
    \textcolor{red}{\textbf{A. 丁立人}} \qquad B. 韦奕 \qquad C. 王玥 \qquad D. 卜祥志 \\
    \textcolor{blue}{\textbf{【现场解说】}：2800分是国际象棋界的“神之领域”，全球历史上仅有极少数顶尖天才达到过此成就！}

    \item \textbf{在“将军三次即获胜 (3-Check)”变体中，如果一次移动同时达成“双反击将军”，算几次将军？} \\
    \textcolor{red}{\textbf{A. 1次}} \qquad B. 2次 \qquad C. 3次 \qquad D. 0次 \\
    \textcolor{blue}{\textbf{【现场解说】}：每走一步棋，无论产生几个将军射线，只计为1次 Check 动作。在这模式里，丢子是小事，被将军才是致命的！}


\item \textbf{中国科大校学生深蓝国际象棋协会正式成立于哪一年？} \\
    \textcolor{red}{\textbf{答案：B. 2003年}} \\
    \textcolor{blue}{\textbf{【社团记忆】：} 咱们协会成立于2003年，至今已有20余年历史，是校内底蕴深厚的智力运动社团。}

    \item \textbf{深蓝协会秉承的宗旨是：“\uline{\hspace{1em}}演绎风云变幻，\uline{\hspace{1em}}提升棋艺思维”。} \\
    \textcolor{red}{\textbf{答案：B. 棋盘/布阵}} \\
    \textcolor{blue}{\textbf{【社团记忆】：} 这是写在社团章程里的理念，强调棋盘上的逻辑推演与策略规划。}

    \item \textbf{2003年10月4日，深蓝协会在校内创办了哪项极具品牌影响力的系列赛事？} \\
    \textcolor{red}{\textbf{答案：C. 深蓝杯}} \\
    \textit{\textbf{【社团记忆】：} “深蓝杯”诞生于建社同年国庆期间，是协会最核心的传统赛事。}

    \item \textbf{在明日（4月19日）的跨校友谊赛中，深蓝协会将与哪所高校的棋队展开线上对决？} \\
    \textcolor{red}{\textbf{答案：C. 浙江大学}} \\
    \textcolor{blue}{\textbf{【实时速递】：} 明天我们将在 Lichess 平台与浙大进行校际联赛，欢迎大家关注并为科大助威！}

    \item \textbf{现代国际象棋的规则（特别是王后和象的强力走法）大约是在哪个时期定型的？} \\
    \textcolor{red}{\textbf{答案：B. 15世纪末的欧洲}} \\
    \textcolor{blue}{\textbf{【历史科普】：} 虽然起源于印度，但在15世纪末的西班牙和意大利，王后才变成了全场最强的棋子。}

    \item \textbf{历史上第一位官方公认的国际象棋世界冠军是？（1886年夺冠）} \\
    \textcolor{red}{\textbf{答案：B. 威廉·斯坦尼茨}} \\
    \textcolor{blue}{\textbf{【历史科普】：} 斯坦尼茨不仅是首位世界冠军，还是现代定位理论与阵地战防守的先驱。}

    \item \textbf{国际象棋联合会（FIDE）于1924年成立于哪个城市？} \\
    \textcolor{red}{\textbf{答案：C. 巴黎}} \\
    \textcolor{blue}{\textbf{【历史科普】：} 国际棋联的座右铭是“Gens una sumus”（我们同属一家）。}

    \item \textbf{中国第一位夺得国际象棋女子世界冠军，打破欧洲多年垄断的棋手是谁？} \\
    \textcolor{red}{\textbf{答案：B. 谢军}} \\
    \textcolor{blue}{\textbf{【中国崛起】：} 1991年谢军夺冠，标志着中国国际象棋“四步走”战略迈出了第一步。}

    \item \textbf{1956年，中国将国际象棋列为正式体育项目，当时大力推动此决策的元帅是？} \\
    \textcolor{red}{\textbf{答案：C. 陈毅}} \\
    \textcolor{blue}{\textbf{【中国崛起】：} 陈毅元帅曾说：“棋盘也是战场”，他极力推动了中国三棋的普及。}

    \item \textbf{浙江哪座城市因为名手辈出，被授予中国首个“国际象棋之城”的美誉？} \\
    \textcolor{red}{\textbf{答案：C. 温州}} \\
    \textcolor{blue}{\textbf{【中国崛起】：} 温州不仅孕育了诸宸，也是现任世界冠军丁立人的故乡。}

    \item \textbf{中国天才少女侯逸凡在多少岁时首次夺魁，成为历史上最年轻的棋后？} \\
    \textcolor{red}{\textbf{答案：B. 16岁}} \\
    \textcolor{blue}{\textbf{【中国崛起】：} 2010年，16岁的侯逸凡夺冠，打破了此前维持了数十年的“最年轻”纪录。}

    \item \textbf{在原子棋（Atomic）规则中，如果你用王直接吃掉对方的棋子，会发生什么？} \\
    \textcolor{red}{\textbf{答案：B. 你的王会自爆，你输了}} \\
    \textcolor{blue}{\textbf{【变体趣味】：} 原子棋中吃子即爆炸，由于王不能自爆，所以王无法进行捕获动作。}

    \item \textbf{在“疯狂乱斗（Crazyhouse）”中，吃掉的棋子可以空降回棋盘，但“兵”不能空降在？} \\
    \textcolor{red}{\textbf{答案：A. 1行和8行}} \\
    \textcolor{blue}{\textbf{【变体趣味】：} 为了平衡性，空降的兵不能直接放置在底线进行升变。}

    \item \textbf{国际象棋历史上，统治世界冠军宝座时间最长的棋手是？（长达27年）} \\
    \textcolor{red}{\textbf{答案：A. 拉斯克}} \\
    \textcolor{blue}{\textbf{【历史科普】：} 拉斯克从1894年一直到1921年都是世界冠军，他本身也是一位杰出的数学家。}

    \item \textbf{1997年击败世界冠军卡斯帕罗夫的超级电脑“深蓝”是由哪家公司研发的？} \\
    \textcolor{red}{\textbf{答案：B. IBM}} \\
    \textcolor{blue}{\textbf{【背景解析】：} 协会以此命名，既是对这段人机博弈史的纪念，也体现了科大学子探索科技与逻辑的追求。}

\item \textbf{深蓝协会的另一项招牌赛事“闪电杯”超快棋赛，首届比赛是在哪一年举行的？} \\
    \textcolor{red}{\textbf{答案：C. 2022年}} \\
    \textit{\textbf{【社团史】：} 第一届“闪电杯”于2022年4月16日在东区体育教学楼201教室落子，标志着协会超快棋传统的开始。}

    \item \textbf{在第二届“闪电杯”中，哪位同学以8战全胜的统治级表现封王，创造了赛事史上首个全胜夺冠纪录？} \\
    \textcolor{red}{\textbf{答案：B. 高鹤宇}} \\
    \textit{\textbf{【社团史】：} 2023年4月1日，高鹤宇同学在中区棋牌室展现了强大的统治力，至今这仍是“闪电杯”的一段佳话。}

    \item \textbf{关于明日（4月19日）与浙大的友谊赛，专业组第一阶段采用的赛制模式是？} \\
    \textcolor{red}{\textbf{答案：C. 团队竞技场 (Team Battle Arena)}} \\
    \textit{\textbf{【活动解析】：} 这种模式会自动避免同校选手互配，确保选手“一致对外”，积分最高的选手将进入决赛。}

    \item \textbf{在浙大 x 中科大友谊赛的决赛阶段，每方选手的对局用时是多少？} \\
    \textcolor{red}{\textbf{答案：C. 10+5}} \\
    \textit{\textbf{【活动解析】：} 决赛阶段为了考验深度计算，用时从预赛的 5+2 增加到了 10+5。}

    \item \textbf{国际象棋中，如果一方的王没有被将军，但己方所有棋子都没有合法走法，这被称为什么？} \\
    \textcolor{red}{\textbf{答案：B. 逼和 (Stalemate)}} \\
    \textit{\textbf{【知识科普】：} 逼和在国际象棋中计为和盘（Draw）。这是劣势方寻求反转的重要战术。}

    \item \textbf{在“双人蝉联象棋 (Bughouse)”变体中，最核心的规则特点是什么？} \\
    \textcolor{red}{\textbf{答案：B. 两人组队且队友吃掉的棋子可以交给你使用}} \\
    \textit{\textbf{【变体科普】：} 这是深蓝协会嘉年华中最受欢迎的项目之一，极度考验队友间的默契与棋子“空降”的时机。}

    \item \textbf{国际棋联 (FIDE) 授予的最高永久头衔是什么？} \\
    \textcolor{red}{\textbf{答案：C. 特级大师 (GM)}} \\
    \textit{\textbf{【知识科普】：} 特级大师称号是棋手的终身荣耀，丁立人、谢军等名将均为 GM 头衔。}

    \item \textbf{当一名兵移动到对方底线时，它可以升变为除了哪种棋子以外的任何兵种？} \\
    \textcolor{red}{\textbf{答案：B. 王}} \\
    \textit{\textbf{【知识科普】：} 兵可以升变为后、车、象、马。在某些特殊残局中，升变为马或车（低升变）反而可能是致胜关键。}

    \item \textbf{在“吃过路兵”规则中，这个动作必须在对方兵前进两格后的第几步内完成？} \\
    \textcolor{red}{\textbf{答案：A. 立即下一步}} \\
    \textit{\textbf{【知识科普】：} “吃过路兵”权利稍纵即逝，如果你在当下没有执行，之后就不能再针对那枚兵使用此规则。}

    \item \textbf{深蓝协会目前正致力于建立数字化档案，主要整理的内容包括哪些？} \\
    \textcolor{red}{\textbf{答案：D. 以上皆是}} \\
    \textit{\textbf{【社团发展】：} 我们正在整理自2003年建社以来的珍贵影像和荣誉，旨在打造一个“有温度、有传承”的深蓝大家庭。}
\item \textbf{答案：C. 国王} \\
    \textit{\textbf{【解析】}：兵在移动到对方底线时可以升变为除了王和兵以外的任何棋子（后、车、马、象）。升变为王后是最常见的，但在某些特定战术下（如避免逼和）可能会升变为马或车。}

    \item \textbf{答案：D. 参与易位的“车”受到攻击} \\
    \textit{\textbf{【解析】}：王车易位的禁忌包括：王被将军时、易位过程中王经过的格子受攻击时、王和车动过时、王易位后的格子受攻击时。但“车”受到攻击或易位过程中“车”经过受攻击的格子是不受限制的。}

    \item \textbf{答案：B. 3个} \\
    \textit{\textbf{【解析】}：通常的分值计算为：兵1分，马3分，象3.25-3.5分（通常简化为3分），车5分，后9分。象因为能远程控制斜线，在开放局面下略强于马。}

    \item \textbf{答案：A. 立即下一回合} \\
    \textit{\textbf{【解析】}：吃过路兵是一种特殊的吃子方式，权力的行使具有时效性，必须在对方兵双格挺进后的紧接着那一回合执行。}

    \item \textbf{答案：B. 1... c5} \\
    \textit{\textbf{【解析】}：西西里防御（Sicilian Defence）是应对1.e4最强有力且最流行的黑方开局，其核心在于非对称的中心争夺。}

    \item \textbf{答案：D. 在没有明确目标时多次移动同一个棋子} \\
    \textit{\textbf{【解析】}：开局讲究效率。多次移动同一个棋子会浪费“时机”（Tempo），导致出子落后，容易给对手留下攻击机会。}

    \item \textbf{答案：A. f7} \\
    \textit{\textbf{【解析】}：f7格（黑方）或f2格（白方）在开局阶段仅由王一人防守，是阵型中最脆弱的点，著名的“学子棋”和“肝油攻击”都聚焦于此。}

    \item \textbf{答案：C. 后翼弃兵} \\
    \textit{\textbf{【解析】}：白方第一步走1.d4被称为封闭或半封闭式开局。后翼弃兵（Queen's Gambit）是1.d4 d5 2.c4。由于《后翼弃兵》同名美剧的火热，这一开局广为人知。}

    \item \textbf{答案：B. 对王（Opposition）} \\
    \textit{\textbf{【解析】}：对王是指两个国王在同一行或列上相对，中间相隔奇数个格子。占据对王的一方可以迫使对方国王让位，从而让自己的兵顺利推进。}

    \item \textbf{答案：A. 王能否追上对方的兵} \\
    \textit{\textbf{【解析】}：如果在对方挺兵时，己方国王能一步踏入由该兵到升变格构成的“虚拟正方形”内，国王就能追上该兵。}

    \item \textbf{答案：D. 将死} \\
    \textit{\textbf{【解析】}：将死（Checkmate）意味着对局结束。而逼和（Stalemate）则是王没被将军但无路可走，结果计为和局。}

    \item \textbf{答案：C. 理论上可以，但需要极高技巧} \\
    \textit{\textbf{【解析】}：马象杀单王是国际象棋中最难的基础残局之一，通常需要在50招之内将对方王驱赶至与象同色的角格才能完成。}

    \item \textbf{答案：C. 古印度} \\
    \textit{\textbf{【解析】}：现代国际象棋起源于大约公元6世纪印度的“恰图兰卡”，意为“四色棋”，后经由波尔斯传入欧洲演变为现代规则。}

    \item \textbf{答案：A. 鲍比·菲舍尔} \\
    \textit{\textbf{【解析】}：1972年，菲舍尔在冰岛雷克雅未克击败斯帕斯基，这是历史上关注度最高的世界冠军赛之一。}

    \item \textbf{答案：B. 深蓝 (Deep Blue)} \\
    \textit{\textbf{【解析】}：1997年“深蓝”击败卡斯帕罗夫，开启了计算机国际象棋的新时代。这也是我们“深蓝国际象棋协会”名称的文化渊源。}
\item \textbf{答案：C. 和棋} \\
    \textit{\textbf{【解析】}：当同一局面（包括棋子位置、轮到谁走、易位权及吃过路兵权）在对局中出现三次时，任何一方均可向裁判提议和棋。这是为了防止无意义的循环走法。}

    \item \textbf{答案：A. c1或c8} \\
    \textit{\textbf{【解析】}：在大易位（长位）中，国王向后翼移动两格，最终落在 c1（白方）或 c8（黑方），而车则跳过国王落在 d1 或 d8。}

    \item \textbf{答案：C. 连续50步内既没有棋子被吃，也没有动过兵} \\
    \textit{\textbf{【解析】}：这是为了防止棋手在无法取胜的残局中无休止地拖延。只要连续50个回合内没有吃子和兵的移动，对局即判为和棋。}

    \item \textbf{答案：A. 合法} \\
    \textit{\textbf{【解析】}：只要“吃过路兵”这一动作能合规地消除被“将军”的状态（即吃掉将军的兵或挡住线路），它就是合法的。}

    \item \textbf{答案：B. c4} \\
    \textit{\textbf{【解析】}：1.e4 e5 2.Nf3 Nc6 3.Bc4 是意大利开局。这个位置的象直指黑方脆弱的 f7 点，极具攻击性。}

    \item \textbf{答案：A. 1... e6} \\
    \textit{\textbf{【解析】}：法兰西防御是典型的半封闭式开局，黑方先筑起 e6-d5 的防御堡垒，准备在反击中挑战白方的中心。}

    \item \textbf{答案：C. 中心} \\
    \textit{\textbf{【解析】}：后翼弃兵（1.d4 d5 2.c4）通过引诱黑方兵离开中心（若黑方吃掉 c4 兵），从而让白方能够通过 e4 占据整个中心地带。}

    \item \textbf{答案：A. 被对方马、象等轻子通过攻击“后”来赚取步数} \\
    \textit{\textbf{【解析】}：后的价值最高，对方只需用价值较低的棋子（如马）威慑后，你就必须移动后，这样对方在出子的同时还消耗了你的行动轮次。}

    \item \textbf{答案：B. 对王（Opposition）} \\
    \textit{\textbf{【解析】}：当两王相对，中间只隔一格时，谁不走棋谁就占据了“对王”优势。在残局中，这常被用来强行推兵升变。}

    \item \textbf{答案：A. 搭桥} \\
    \textit{\textbf{【解析】}：在车兵残局中，当自己的王被对方车在底线封锁时，通过将自己的车移动到第4行或第5行来挡住对方车的将军，这个战术生动地称为“搭桥（Building a bridge）”。}

    \item \textbf{答案：C. 车兵残局} \\
    \textit{\textbf{【解析】}：菲利多尔位置是防守方在车兵残局中谋求和棋的最经典方法，通过在第三横线切断对方国王的去路来实现。}

    \item \textbf{答案：C. 马} \\
    \textit{\textbf{【解析】}：由于马的走法是“L”型且可以越过其他棋子，它经常能同时攻击对方两个重要目标（如同时将军并攻击后），让对方防不胜分。}

    \item \textbf{答案：C. 瑞士} \\
    \textit{\textbf{【解析】}：国际棋联（FIDE）目前的总部位于瑞士洛桑，那里也是许多国际体育组织（如国际奥委会）的所在地。}

    \item \textbf{答案：B. 卡帕布兰卡} \\
    \textit{\textbf{【解析】}：何塞·劳尔·卡帕布兰卡以几乎不出错的简明风格著称，他在1916年到1924年间甚至没有输掉过任何一场正式对局。}

    \item \textbf{答案：A. 卡斯帕罗夫} \\
    \textit{\textbf{【解析】}：加里·卡斯帕罗夫是国际象棋史上最伟大的棋手之一。1997年他在纽约输给了升级后的“深蓝”，这被视为人工智能史上的里程碑。}

    \item \textbf{答案：B. 马格努斯·卡尔森} \\
    \textit{\textbf{【解析】}：来自挪威的卡尔森自2011年起长期占据世界第一，他不仅在常规赛表现卓越，在快棋和超快棋领域也具有统治力。}
\end{enumerate}

\vspace{2em}
\begin{center}
    \textbf{祝“薪火传承”校园活动圆满成功！}
\end{center}

\end{document}